\section{The geometric method: a usable arithmetic determinant reduction}
\label{sec:root-method}

There are two distinct multiplicities in the prompt.  An automorphic
restriction multiplicity is the dimension of a space of equivariant
functionals.  A Stepanov multiplicity is an order of vanishing or an
integral divisibility order.  A theorem relating them must include an
integral lattice, a valuation, and compatible maps; a complex
representation-theoretic dimension alone does not supply those data.
The following finite-dimensional form isolates the actual numerical
inequality needed by the arithmetic determinant method.

\paragraph{Finding BR-DET-01 (divisibility versus archimedean size).}
Let $K$ be a number field with ring of integers $\mathcal O_K$.
Let $A$ be an $r\times r$ matrix over $\mathcal O_K$.  Suppose its
$j$th column has every entry in a nonzero integral ideal $\mathfrak a_j$,
and for each embedding $\sigma:K\hookrightarrow\mathbb C$ the
Euclidean norm of that column is at most $C_{\sigma,j}>0$.
Count all embeddings, including both members of a complex pair.
If
\[
 \sum_{j=1}^r\log N\mathfrak a_j
   >\sum_{\sigma:K\hookrightarrow\mathbb C}
                     \sum_{j=1}^r\log C_{\sigma,j},             \tag{D.1}
\]
then $\det A=0$.  If $\det A\ne0$, the reverse weak inequality
is necessary.

\begin{proof}
Expansion by permutations shows
$\det A\in\mathfrak a_1\cdots\mathfrak a_r$.
For $a=\det A\ne0$, inclusion of integral ideals
$(a)\subseteq\mathfrak a_1\cdots\mathfrak a_r$ gives
\[
 \prod_jN\mathfrak a_j
 \leq N(a)=|N_{K/\mathbb Q}(a)|
 =\prod_\sigma|\sigma(a)|.
\]
Hadamard's determinant inequality, applied separately at every
embedding, gives
$|\sigma(a)|\leq\prod_j C_{\sigma,j}$.
Multiplying and taking logarithms proves the reverse weak
inequality.  Its contradiction with (D.1) proves vanishing.
\end{proof}

\paragraph{Finding BR-DET-02 (the exact accumulation law for iterated lifts).}
Let $L_j$ be free finite-rank $\mathcal O_K$-lattices with Hermitian
norms at the embeddings, and let
$\theta_j:L_{j-1}\to L_j$ be $\mathcal O_K$-linear maps.
Suppose
\[
 \theta_j(L_{j-1})\subseteq\mathfrak b_jL_j,
 \qquad \|\theta_{j,\sigma}\|\leq c_{j,\sigma}.
\]
Then
\[
 \theta_m\cdots\theta_1(L_0)
   \subseteq(\mathfrak b_1\cdots\mathfrak b_m)L_m,
 \qquad
 \|\theta_{m,\sigma}\cdots\theta_{1,\sigma}\|
   \leq\prod_{j=1}^m c_{j,\sigma}.                            \tag{D.2}
\]
After an integral evaluation map $e:L_m\to\mathcal O_K^r$, columns
obtained from chosen vectors $v_k\in L_0$ obey BR-DET-01 with
the bound
\[
 C_{\sigma,k}=\|e_\sigma\|
       \left(\prod_{j=1}^m c_{j,\sigma}\right)\|v_{k,\sigma}\|,
\]
or any strictly positive upper bound if a column is zero.
Thus a valid application of ``one order
per lift'' requires the ideal-inclusion assertion in (D.2), and
a valid use of arbitrarily many lifts requires their accumulated
logarithmic gain to exceed their accumulated norm cost.

\begin{proof}
For any ideal $\mathfrak b$ and any $\mathcal O_K$-linear map
$\theta$, one has $\theta(\mathfrak bL)=\mathfrak b\theta(L)$:
both sides consist of finite sums of coefficients from
$\mathfrak b$ times images of elements of $L$.
Induction proves the lattice inclusion.  Submultiplicativity of
operator norms proves the other inequality.  An integral
evaluation map preserves the ideal inclusions; applying the
composite operator to $v_k$ bounds its norm by the product of
the operator norms and $\|v_{k,\sigma}\|$. The preceding
finding applies to the resulting matrix.
\end{proof}

These findings are elementary reductions, not a replacement for
Alp\"oge's arithmetic estimates.  They explain exactly which
content the four-page method would have to export in order to
transfer to another motivic family.  CM-period divisibility and
the integral theta structures supply the ideals.  A Rallis
formula and period estimates may control the archimedean norms.
The Iwasawa averaging and relative trace formula must preserve
the necessary independence and integrality while improving the
cost.  The order-$m$ cancellation in the sketch matters because
it changes the sign of the strict inequality (D.1).  Neither
that cancellation nor a uniform tower estimate is a formal
consequence of the five exported RH declarations.

The exact external comparison points are Gan--Gross--Prasad
(2009), Conjectures 17.1 and 17.3 (local distinguished members),
Theorem 14.1 and the reductions in Theorems 15.1 and 16.1
(multiplicity-one statements in their stated cases), and
Conjecture 24.1 (global periods and central values).  These are
the statuses and numbering of the supplied work; we do not call
all of its conjectures theorems from its presence in the ZIP.
Gan--Qiu--Takeda (2014), Theorem 11.4, is a precise Rallis
inner-product formula, with range, cuspidality, and local
conditions.  Applying it repeatedly requires checking those
conditions at every step.  Its special values generally belong
to higher-rank automorphic $L$-functions, outside the stated
Dirichlet scope.  GRH also allows zeros at the centre; it does
not assert that a central period is nonzero.

\subsection{What this does for crystals and for towers}

For a rank-one abelian Artin--Tate object unramified at a
realization prime $p$, its crystalline realization has Frobenius
eigenvalue $p^{-n}$ times the unramified finite character, in the
convention of BR-AT-01.  Its Newton slope is $-n$ and its
Hodge--Tate weight is $-n$.  The period $t$ with $\varphi(t)=pt$ gives this normalization
(Berger, 2004, \S II.3.1 and II.4.3); an unramified finite twist
changes the eigenvalue only by a root of unity.  This local rank-one
compatibility is already known; GRH is a global assertion about the continued
Euler product.  For a ramified finite character one must first
state a potentially crystalline or other appropriate local
realization, rather than silently use the unramified formula.
For the CM curve $y^2=x^3+1$, the separate argument below does
give a genuine new use of Dirichlet GRH: it controls the supply
of inert, trace-zero primes used by the divisibility mechanism.
Its two-dimensional $H^1$ remains outside the diagonal
Artin--Tate Hodge types.

In a cyclotomic tower of conductor $p^a$, the finite Frobenius
recognition theorem below applies at every fixed level, with
$Q=\operatorname{lcm}(p^a,4)$ and bound
$(\varphi(Q)\log Q)^2$.  This is an explicit effective
consequence for the tower's finite abelian layers.  An inverse
limit of integral Selmer groups, a statement about its
characteristic ideal, or a Vandiver-type vanishing assertion
requires additional integral information.  BR-MIX-01 supplies
an explicit example of information lost on passing to the
weight-graded $L$-function; it prevents substituting a
semisimple spectral calculation for that missing integral
calculation.

\subsection{Arithmetic duality is a precise analogy, not a missing construction}

Artin--Verdier duality is a pairing for constructible torsion
\'etale coefficients on number rings, with compact support and
the real-place conventions.  A standard exact source is Milne
(2006), \emph{Arithmetic Duality Theorems}, Chapter II,
Theorem 3.1(b), p.~177.  Its degree-three duality explains the arithmetic
three-dimensional analogy.  It does not equip
$\operatorname{Spec}\mathbb Z$ with a positive Hilbert-space
cohomology whose determinant is $\xi$.  In fact the infinitely
many zeros of $\xi$ already exclude an ordinary characteristic
polynomial of a single finite-dimensional complex operator.
If such a finite-dimensional operator $A$ existed, its
determinant $\det(sI-A)$ would be a polynomial with at most
$\dim A$ zeros, contradicting the known zero count.
An infinite-dimensional space, a trace class transformation
of the variable as in BR-TR-01, or an appropriate regularized
determinant is additional structure.  This is the precise
analytic point at which the arithmetic-duality analogy and a
Hilbert--P\'olya or Deninger construction must be connected by
a new theorem.

